%% ====================================================================
\section{Desarrollos teóricos integrados}
\label{sec:desarrollos-teoricos-integrados}
\label{sec:anexo-teorico-extendido}
%% ====================================================================

Esta sección integra las ecuaciones avanzadas que completan el marco teórico del
Capítulo~\ref{sec:marco}. El objetivo es mostrar cómo la asignación por teoría de
juegos, la capacidad mecánica efectiva, la seguridad, la batería, la predicción de
densidad y el transporte cooperativo pueden escribirse en una arquitectura matemática
común. Los resultados se presentan con su hipótesis de validez: algunos quedan
demostrados dentro del modelo fluido, otros requieren regularidad, conectividad o
acotación de errores para aplicarse al sistema discreto de AGVs.

\subsection{Lift general por contribución efectiva}
\label{subsec:lift-contribucion}

El modelo homogéneo, la ocupación efectiva, la capacidad heterogénea y la ruta eikonal
son casos particulares de una construcción más general: toda física relevante entra al
juego como una contribución efectiva no negativa.

\begin{definicion}[Contribución efectiva general]
Sea $\xi_i$ el estado interno del robot $i$, $\eta_k$ el descriptor de la tarea $k$,
$\mathcal{M}$ el mapa o modelo compartido y $t$ el tiempo. Una contribución efectiva es
una función
\begin{equation}
  s_{ik}=s_{ik}(\xi_i,\eta_k,\mathcal{M},t)\ge0.
  \label{eq:anexo_s_general}
\end{equation}
El agregado de servicio de la tarea $k$ es
\begin{equation}
  \Zeff_k=\sum_i y_{ik}s_{ik}.
  \label{eq:anexo_Z_general}
\end{equation}
\end{definicion}

\begin{figure}[H]
\centering
\begin{tikzpicture}[
  block/.style={draw=VIUDark,fill=VIULight,rounded corners=2pt,
    align=center,minimum width=3.0cm,minimum height=0.85cm,font=\small},
  arrow/.style={-{Latex[length=2mm]},thick,color=VIUOrange!90!black}
]
\node[block] (physics) {física local\\$c_i,T_k,B_i,\rho(x)$};
\node[block, right=0.8cm of physics] (s) {contribución\\$s_{ik}$};
\node[block, right=0.8cm of s] (Z) {agregado\\$\Zeff_k=\sum_i y_{ik}s_{ik}$};
\node[block, below=0.85cm of Z] (F) {payoff\\$F_{ik}=\rho_k\Phi_k(\Zeff_k)s_{ik}$};
\node[block, below=0.85cm of s] (Smith) {Smith\\$\dot y_i$};
\node[block, below=0.85cm of physics] (motion) {movimiento\\$p_i,\theta_i$};
\draw[arrow] (physics) -- (s);
\draw[arrow] (s) -- (Z);
\draw[arrow] (Z) -- (F);
\draw[arrow] (F) -- (Smith);
\draw[arrow] (Smith) -- (motion);
\draw[arrow] (motion) -- (physics);
\end{tikzpicture}
\caption{Lift por contribución efectiva. La física no entra como regla externa, sino
como unidad de servicio $s_{ik}$ que preserva la estructura de potencial.}
\label{fig:anexo_lift}
\end{figure}

\begin{teorema}[Lift preservador de potencial por contribución efectiva]
\textbf{[ADOPTADO]}
Si $s_{ik}$ no depende de $y_i$ y el payoff se define como
\begin{equation}
  F_{ik}=\rho_k\Phi_k(\Zeff_k)s_{ik},
  \qquad
  F_{i0}=\rho_0,
  \label{eq:anexo_payoff_lift}
\end{equation}
entonces el potencial
\begin{equation}
  V(y,\xi)=
  \rho_0\sum_i y_{i0}
  +
  \sum_k\rho_kG_k(\Zeff_k),
  \qquad
  G_k'(z)=\Phi_k(z),
  \label{eq:anexo_potencial_lift}
\end{equation}
satisface
\begin{equation}
  \frac{\partial V}{\partial y_{ik}}=F_{ik}.
  \label{eq:anexo_grad_lift}
\end{equation}
\end{teorema}

\begin{proof}
Por regla de la cadena,
\[
  \frac{\partial V}{\partial y_{ik}}
  =
  \rho_kG_k'(\Zeff_k)\frac{\partial \Zeff_k}{\partial y_{ik}}
  =
  \rho_k\Phi_k(\Zeff_k)s_{ik}.
\]
La última expresión coincide con \eqref{eq:anexo_payoff_lift}. El pago de inactividad
se obtiene directamente de $\partial(\rho_0\sum_i y_{i0})/\partial y_{i0}=\rho_0$.
\end{proof}

El teorema se interpreta con el estado físico $\xi$ congelado al derivar respecto de
$y$. Si $s_{ik}$ depende directamente de $y$, aparecen términos adicionales
$\partial s_{jk}/\partial y_{ik}$ y la condición de potencial debe reanalizarse.

El caso homogéneo se recupera con $s_{ik}=1$; la ocupación efectiva con
$s_{ik}=\Psi(d_{ik})$; la capacidad heterogénea con $s_{ik}=c_i$; y el caso eikonal con
$s_{ik}=c_i\Psi(T_k(p_i))$.

\subsection{Condición inversa de integrabilidad}
\label{subsec:integrabilidad-inversa}

El lift anterior no es sólo una elección elegante. También aparece como condición
necesaria para que el juego conserve un potencial exacto cuando el agregado y el payoff
usan factores físicos.

\begin{proposicion}[Condición inversa de integrabilidad]
\textbf{[ADOPTADO]}
Sea
\begin{equation}
  Z_k=\sum_i a_{ik}y_{ik},
  \qquad
  F_{ik}=\rho_k\Phi_k(Z_k)b_{ik},
  \label{eq:anexo_integrabilidad_setup}
\end{equation}
con $a_{ik}>0$ y $b_{ik}>0$. Para que exista un potencial exacto separable por tarea,
las derivadas cruzadas deben satisfacer
\begin{equation}
  \frac{\partial F_{ik}}{\partial y_{jk}}
  =
  \frac{\partial F_{jk}}{\partial y_{ik}},
  \qquad i\ne j.
  \label{eq:anexo_derivadas_cruzadas}
\end{equation}
Si $\rho_k\Phi_k'(Z_k)\ne0$, esta condición equivale a
\begin{equation}
  a_{jk}b_{ik}=a_{ik}b_{jk}
  \quad\Longleftrightarrow\quad
  b_{ik}=\alpha_k a_{ik}
  \label{eq:anexo_condicion_integrabilidad}
\end{equation}
para alguna escala $\alpha_k>0$ dependiente sólo de la tarea.
\end{proposicion}

\begin{proof}
Derivando \eqref{eq:anexo_integrabilidad_setup},
\[
  \frac{\partial F_{ik}}{\partial y_{jk}}
  =
  \rho_k\Phi_k'(Z_k)a_{jk}b_{ik},
  \qquad
  \frac{\partial F_{jk}}{\partial y_{ik}}
  =
  \rho_k\Phi_k'(Z_k)a_{ik}b_{jk}.
\]
Si el factor común no se anula, la igualdad de derivadas cruzadas equivale a
$a_{jk}b_{ik}=a_{ik}b_{jk}$. Dividiendo por $a_{ik}a_{jk}$ se obtiene que
$b_{ik}/a_{ik}=b_{jk}/a_{jk}$ para todo par $i,j$, por lo que existe una escala
$\alpha_k$ con $b_{ik}=\alpha_ka_{ik}$.
\end{proof}

\begin{observacion}
Esta proposición formaliza por qué la ocupación efectiva no es una corrección estética:
el mismo factor físico que entra al agregado debe entrar al payoff marginal. Si el
agregado usa $a_{ik}=\Psi_{ik}$ y el payoff usa $b_{ik}=1$, la simetría cruzada falla.
La equivalencia se lee por componentes activas: si algún factor se anula, o si
$\Phi_k'(Z_k)=0$ en una zona plana idealizada, la condición debe evaluarse sólo en las
direcciones donde la derivada cruzada es informativa.
\end{observacion}

\subsection{Oferta heterogénea por tramos}
\label{subsec:oferta-tramos}

La función $h(Q)$ del Capítulo~\ref{sec:marco} puede escribirse de forma explícita. Sea
\[
  C_p=c^pN^p,
  \qquad
  A_p=\sum_{q=p}^{P}C_q,
  \qquad
  A_{P+1}=0,
\]
donde $A_p$ es la capacidad acumulada desde la clase $p$ hasta la clase más pesada.
Se asume el orden $0<c^1<\cdots<c^P$ y una relajación fluida de masa
$a_k^p\in\mathbb{R}_+$. Con robots enteros, $h(Q)$ deja de ser la curva lineal por
tramos y se obtiene una función escalonada asociada a un problema de selección discreta.
Como los robots pesados se despliegan primero, si $Q\in[A_{p+1},A_p]$ entonces
\begin{equation}
  h(Q)=
  \sum_{q=p+1}^{P}N^q
  +
  \frac{Q-A_{p+1}}{c^p}.
  \label{eq:anexo_h_por_tramos}
\end{equation}
La pendiente del tramo es
\begin{equation}
  h'(Q)=\frac{1}{c^p}.
  \label{eq:anexo_pendiente_h}
\end{equation}

\begin{figure}[H]
\centering
\begin{tikzpicture}[scale=0.96]
  \draw[-{Latex[length=2mm]},thick] (0,0) -- (8.2,0) node[right] {$Q$};
  \draw[-{Latex[length=2mm]},thick] (0,0) -- (0,4.7) node[above] {$h(Q)$};
  \draw[VIUOrange,very thick] (0.6,0.5) -- (2.4,1.0) -- (4.8,2.15) -- (7.4,4.0);
  \draw[dashed,VIUGray] (2.4,0) -- (2.4,1.0);
  \draw[dashed,VIUGray] (4.8,0) -- (4.8,2.15);
  \node[anchor=north,font=\small] at (2.4,0) {$A_{p+1}$};
  \node[anchor=north,font=\small] at (4.8,0) {$A_p$};
  \node[font=\small,VIUDark] at (3.65,1.85) {pendiente $1/c^p$};
  \draw[decorate,decoration={brace,amplitude=5pt},VIUGray]
    (2.4,-0.35) -- (4.8,-0.35)
    node[midway,below=6pt,font=\small] {$Q\in[A_{p+1},A_p]$};
\end{tikzpicture}
\caption{Forma explícita de $h(Q)$ por tramos. Cada tramo corresponde a la clase de
capacidad marginal que todavía se está desplegando.}
\label{fig:anexo_h_tramos}
\end{figure}

Con $\rho_p=\rho_0/c^p$, el orden de capacidades implica
\[
  \rho_P<\rho_{P-1}<\cdots<\rho_1.
\]
Para respetar el dominio del logaritmo, se define
\[
  W_k^M(\lambda)=
  \begin{cases}
  0,
  & \lambda\ge \rho_k\Phi_k(0),\\[4pt]
  M_k+\dfrac{1}{\beta}\log\!\left(\dfrac{\rho_k}{\lambda}-1\right),
  & 0<\lambda<\rho_k\Phi_k(0),
  \end{cases}
\]
y la demanda agregada
\[
  D(\lambda)=\sum_k W_k^M(\lambda),
\]
los tres regímenes se escriben así:

\begin{description}
  \item[Tramo.] Si
  \begin{equation}
    A_{p+1}<D(\rho_p)<A_p,
    \label{eq:anexo_condicion_tramo}
  \end{equation}
  entonces $\lambda^\star=\rho_p$.

  \item[Kink.] Si
  \begin{equation}
    D(\rho_p)\ge A_p\ge D(\rho_{p-1}),
    \label{eq:anexo_condicion_kink}
  \end{equation}
  entonces
  \begin{equation}
    Q^\star=A_p,
    \qquad
    D(\lambda^\star)=A_p,
    \qquad
    \lambda^\star\in[\rho_p,\rho_{p-1}].
    \label{eq:anexo_kink}
  \end{equation}

  \item[Escasez total.] Si
  \begin{equation}
    D(\rho_1)>A_1,
    \label{eq:anexo_condicion_escasez}
  \end{equation}
  entonces
  \begin{equation}
    Q^\star=A_1,
    \qquad
    D(\lambda^\star)=A_1,
    \qquad
    \lambda^\star\ge\rho_1.
    \label{eq:anexo_escasez}
  \end{equation}
\end{description}

\subsection{ISS práctica del juego perturbado}
\label{subsec:iss-practica}

El sistema implementado no recibe payoffs exactos. Usa estimaciones locales, mapas
aproximados, discretización y ruido. Es útil escribir la robustez como una cota ISS.

\begin{proposicion}[ISS práctica por error de payoff]
\textbf{[CONDICIONADO]}
Sea
\[
  \widehat F(t)=F(t)+e(t),
  \qquad
  \norm{e(t)}_\infty\le\bar e,
\]
y sea $W(y)=V^\star-V(y)$ una función de error de potencial respecto al conjunto de
Nash $\mathcal{N}$. Supóngase que el sistema nominal cumple localmente
\begin{equation}
  \dot W_{\mathrm{nom}}
  \le
  -a\,d(y,\mathcal{N})^2,
  \qquad a>0.
  \label{eq:anexo_iss_nominal}
\end{equation}
Si la perturbación entra como
\begin{equation}
  \dot W
  \le
  -a\,d(y,\mathcal{N})^2
  +
  b\,d(y,\mathcal{N})\bar e,
  \label{eq:anexo_iss_perturbada}
\end{equation}
entonces
\begin{equation}
  \dot W
  \le
  -\frac{a}{2}d(y,\mathcal{N})^2
  +
  \frac{b^2}{2a}\bar e^2,
  \label{eq:anexo_iss_young}
\end{equation}
y, bajo equivalencia local entre $W$ y $d(y,\mathcal{N})^2$,
\begin{equation}
  \limsup_{t\to\infty}d(y(t),\mathcal{N})
  \le
  C'\bar e.
  \label{eq:anexo_iss_limsup}
\end{equation}
\end{proposicion}

\begin{proof}
Aplicando Young al término cruzado,
\[
  b\,d\,\bar e
  \le
  \frac{a}{2}d^2+\frac{b^2}{2a}\bar e^2.
\]
Sustituir en \eqref{eq:anexo_iss_perturbada} da
\eqref{eq:anexo_iss_young}. La cota de limsup se obtiene con el argumento estándar de
estabilidad práctica: fuera de una vecindad proporcional a $\bar e$, el término
negativo domina.
\end{proof}

La cota abstracta se conecta con errores medibles mediante
\[
  \bar e
  =
  \bar e_{\mathrm{DAC}}
  +
  \bar e_{\Psi}
  +
  \bar e_{\mathrm{disc}}
  +
  \bar e_{\mathrm{hyb}}.
\]
Para el canal DAC homogéneo,
\[
  \bar e_{\mathrm{DAC}}
  \le
  r_{\max}\frac{\beta}{4}\bar\varepsilon,
\]
y en el lift heterogéneo, si $s_{ik}\le c_{\max}$ y el error agregado es
$\bar\varepsilon_Z$,
\[
  \bar e_{\mathrm{DAC}}
  \le
  c_{\max}r_{\max}\frac{\beta}{4}\bar\varepsilon_Z.
\]

\subsection{ORCA económico completo}
\label{subsec:orca-economico}

La evasión local de colisiones debe respetar la criticidad de cada robot para la
coalición. La responsabilidad económica reparte la corrección relativa de ORCA usando
pesos derivados de pérdida de potencial.

\begin{proposicion}[Responsabilidad económica par-a-par]
\textbf{[CONDICIONADO]}
Sea $u_{ij}$ la corrección relativa mínima para que el par $(i,j)$ sea seguro. Defínase
\begin{equation}
  \alpha_{ij}=\frac{w_j}{w_i+w_j},
  \qquad
  \alpha_{ji}=\frac{w_i}{w_i+w_j},
  \qquad
  \alpha_{ij}+\alpha_{ji}=1.
  \label{eq:anexo_alpha_orca}
\end{equation}
Las velocidades corregidas son
\begin{equation}
  v_i^{\mathrm{new}}=v_i^{\mathrm{des}}+\alpha_{ij}u_{ij},
  \qquad
  v_j^{\mathrm{new}}=v_j^{\mathrm{des}}-\alpha_{ji}u_{ij}.
  \label{eq:anexo_vel_orca}
\end{equation}
Entonces
\begin{equation}
  v_i^{\mathrm{new}}-v_j^{\mathrm{new}}
  =
  v_i^{\mathrm{des}}-v_j^{\mathrm{des}}+u_{ij},
  \label{eq:anexo_rel_orca}
\end{equation}
por lo que la corrección relativa completa se conserva.
\end{proposicion}

\begin{proof}
Restando las dos ecuaciones de \eqref{eq:anexo_vel_orca},
\[
  v_i^{\mathrm{new}}-v_j^{\mathrm{new}}
  =
  v_i^{\mathrm{des}}-v_j^{\mathrm{des}}
  +(\alpha_{ij}+\alpha_{ji})u_{ij}.
\]
Usando \eqref{eq:anexo_alpha_orca}, el factor entre paréntesis es $1$.
\end{proof}

En implementación conviene recortar la responsabilidad para evitar que un robot absorba
toda la corrección por una diferencia extrema de pesos:
\begin{equation}
  \widetilde\alpha_{ij}
  =
  \clip\!\left(
  \frac{w_j}{w_i+w_j},
  \alpha_{\min},
  1-\alpha_{\min}
  \right),
  \qquad
  \widetilde\alpha_{ji}=1-\widetilde\alpha_{ij}.
  \label{eq:anexo_alpha_orca_clip}
\end{equation}

El peso recomendado no mide tentación instantánea, sino pérdida finita de potencial:
\begin{equation}
  w_i=
  \varepsilon+
  \sum_k
  y_{ik}\rho_k
  \left[
    G_k(\Zeff_k)-G_k(\Zeff_k-s_{ik})
  \right]_+.
  \label{eq:anexo_peso_orca}
\end{equation}

\begin{observacion}[Por qué un peso escalar dentro del QP no basta]
Si $w_i>0$, entonces
\begin{equation}
  \argmin_{v\in C}w_i\norm{v-v_d}^2
  =
  \argmin_{v\in C}\norm{v-v_d}^2.
  \label{eq:anexo_bug_peso_qp}
\end{equation}
Multiplicar la función objetivo por un escalar positivo no cambia el minimizador. Por
eso el peso debe entrar en el reparto de responsabilidad $\alpha_{ij}$, no como simple
factor de coste en un QP individual.
\end{observacion}

\begin{observacion}[Factibilidad multiagente]
La garantía par-a-par no implica factibilidad global si la intersección de todos los
semiplanos ORCA del robot está vacía. En ese caso se requiere una relajación penalizada,
prioridad de parada segura o reducción temporal de velocidad.
\end{observacion}

\begin{figure}[H]
\centering
\begin{tikzpicture}[
  robot/.style={circle,draw=VIUDark,fill=VIULight,minimum size=0.75cm,font=\small},
  arrow/.style={-{Latex[length=2mm]},thick},
  lab/.style={font=\small,align=center}
]
\node[robot] (i) at (0,0) {$i$};
\node[robot] (j) at (5,0) {$j$};
\draw[arrow,VIUOrange] (i) -- ++(1.7,0.65) node[midway,above,lab] {$\alpha_{ij}u_{ij}$};
\draw[arrow,VIUOrange] (j) -- ++(-1.2,-0.65) node[midway,below,lab] {$-\alpha_{ji}u_{ij}$};
\draw[arrow,VIUGray,dashed] (i) -- ++(1.2,0) node[midway,below,lab] {$v_i^{des}$};
\draw[arrow,VIUGray,dashed] (j) -- ++(-1.2,0) node[midway,above,lab] {$v_j^{des}$};
\node[lab] at (2.5,1.35) {$w_i$ alto $\Rightarrow$ $i$ desvía menos};
\node[lab] at (2.5,-1.25) {$\alpha_{ij}+\alpha_{ji}=1$ conserva la corrección relativa};
\end{tikzpicture}
\caption{ORCA económico: el reparto de la corrección depende de la pérdida de potencial
que causaría desviar a cada robot.}
\label{fig:anexo_orca}
\end{figure}

\subsection{Retornabilidad energética y relevo dinámico}
\label{subsec:retornabilidad-handover}

La batería por sí sola no mide seguridad. Un robot puede tener carga suficiente para
seguir empujando pero no para regresar al cargador. La variable robusta es la
retornabilidad:
\begin{equation}
  h_i(p_i,B_i)=
  B_i-B_{\mathrm{crit}}-E_{\mathrm{return}}(p_i).
  \label{eq:anexo_retornabilidad}
\end{equation}
La recarga puede entrar como una tarea de supervivencia con payoff acotado:
\begin{equation}
  F_{i,\mathrm{ch}}
  =
  \rho_{\mathrm{ch}}
  \sigma(h_{\mathrm{safe}}-h_i)
  \Psi(T_{\mathrm{ch}}(p_i)),
  \label{eq:anexo_payoff_recarga_sigmoidal}
\end{equation}
o como una variante saturada de barrera:
\begin{equation}
  F_{i,\mathrm{ch}}
  =
  \sat\!\left(\frac{\omega}{h_i(p_i,B_i)}\right)
  \Psi(T_{\mathrm{ch}}(p_i)).
  \label{eq:anexo_payoff_recarga_barrera}
\end{equation}

\begin{proposicion}[Relevo dinámico inducido por déficit anticipado]
\textbf{[HIPÓTESIS EXTENDIDA]}
Supóngase que un robot $A$ debe salir a recarga y que
\[
  F_{A,\mathrm{ch}}>F_{A,k}.
  \label{eq:anexo_A_sale}
\]
Entonces Smith induce localmente
\[
  \dot y_{A,\mathrm{ch}}>0,
  \qquad
  \dot y_{A,k}<0.
  \label{eq:anexo_smith_recarga}
\]
Si se predice
\begin{equation}
  \widehat{\Zeff}_k(t+\tau)
  =
  \Zeff_k(t)+\tau\dot{\Zeff}_k(t)
  \downarrow,
  \label{eq:anexo_deficit_anticipado}
\end{equation}
entonces $\Phi_k(\widehat{\Zeff}_k)$ aumenta. Si existe otro robot $B$ con margen
\begin{equation}
  F_{B,k}^{\tau}>F_{B,\ell}^{\tau}+\eta,
  \qquad \eta>0,
  \label{eq:anexo_B_entra}
\end{equation}
entonces Smith transfiere masa de $B$ hacia la tarea $k$. La salida anticipada de $A$
abre déficit y recluta a $B$ sin una regla explícita de relevo.
\end{proposicion}

\begin{proof}
La primera parte sigue directamente de los términos positivos de Smith: si recarga paga
más que la tarea $k$, hay flujo de masa desde $k$ hacia recarga. La salida reduce
$\dot{\Zeff}_k$ por la contribución efectiva perdida, lo que aumenta la sigmoide
decreciente de déficit. El margen \eqref{eq:anexo_B_entra} activa el mismo argumento de
Smith para el robot $B$.
\end{proof}

El relevo no queda garantizado por la batería sola. Es un relevo inducido bajo condición
de margen: si no existe un robot candidato que satisfaga \eqref{eq:anexo_B_entra}, la
salida energética puede reducir el servicio sin reemplazo inmediato.

\begin{figure}[H]
\centering
\begin{tikzpicture}[
  block/.style={draw=VIUDark,fill=VIULight,rounded corners=2pt,
    align=center,minimum width=2.75cm,minimum height=0.85cm,font=\small},
  arrow/.style={-{Latex[length=2mm]},thick,color=VIUOrange!90!black}
]
\node[block] (battery) {$h_A$ bajo\\recarga domina};
\node[block, right=0.7cm of battery] (leave) {$A$ sale\\$\dot y_{A,k}<0$};
\node[block, right=0.7cm of leave] (def) {déficit futuro\\$\widehat{\Zeff}_k\downarrow$};
\node[block, below=0.8cm of def] (price) {payoff sube\\$\Phi_k\uparrow$};
\node[block, left=0.7cm of price] (enter) {$B$ entra\\$\dot y_{B,k}>0$};
\node[block, left=0.7cm of enter] (load) {quórum\\se preserva};
\draw[arrow] (battery)--(leave);
\draw[arrow] (leave)--(def);
\draw[arrow] (def)--(price);
\draw[arrow] (price)--(enter);
\draw[arrow] (enter)--(load);
\end{tikzpicture}
\caption{Relevo dinámico: la salida energética de un robot crea una señal de déficit
que recluta reemplazo mediante el mismo payoff.}
\label{fig:anexo_handover}
\end{figure}

\subsection{Smith anticipatorio en tiempo muestreado}
\label{subsec:smith-sampled-data}

La forma continua $F^\tau=F+\tau\dot F$ puede introducir un bucle algebraico porque
$\dot F$ depende de $\dot y$. Una implementación muestreada evita ese bucle usando la
derivada anterior:
\begin{equation}
  F_m^\tau
  =
  \nabla V(y_m)
  +
  \tau\nabla^2V(y_m)\dot y_{m-1}
  +
  e_m.
  \label{eq:anexo_smith_sampled}
\end{equation}
El error por muestreo es
\begin{equation}
  e_m^\Delta
  =
  \tau\nabla^2V(y_m)(\dot y_{m-1}-\dot y_m).
  \label{eq:anexo_error_sampled}
\end{equation}
Si
\[
  \norm{\dot y_m-\dot y_{m-1}}\le L_\Delta T_s,
\]
entonces
\begin{equation}
  \norm{e_m^\Delta}
  \le
  \tau\norm{\nabla^2V(y_m)}L_\Delta T_s.
  \label{eq:anexo_cota_error_sampled}
\end{equation}

\begin{proposicion}[Cota candidata de Smith anticipatorio muestreado]
\textbf{[HIPÓTESIS EXTENDIDA]}
Bajo suavidad local y error acotado, una cota de descenso práctico para un paso $h$ toma
la forma
\begin{equation}
  W_{m+1}-W_m
  \le
  -h c_\tau\norm{\dot y_m}^2
  +
  hC(\bar e^2+T_s^2+h^2).
  \label{eq:anexo_descenso_sampled}
\end{equation}
En consecuencia,
\begin{equation}
  \limsup_{m\to\infty}d(y_m,\mathcal{N})
  \le
  K(\bar e+T_s+h).
  \label{eq:anexo_limsup_sampled}
\end{equation}
\end{proposicion}

\begin{observacion}[Condición implícita de contracción]
Si se intenta resolver la forma implícita
\[
  \dot y=S\left(y,\nabla V+\tau\nabla^2V\,\dot y\right),
\]
con el operador
\[
  \mathcal{T}(v)=S\left(y,\nabla V+\tau\nabla^2V\,v\right),
\]
una condición tipo Banach es
\begin{equation}
  \tau L_S\norm{\nabla^2V}<1.
  \label{eq:anexo_contraccion_implicita}
\end{equation}
Esta es la condición correcta para hablar de contracción del bucle algebraico.
\end{observacion}

El término anticipatorio debe leerse como amortiguamiento en coordenadas agregadas:
si una dirección microscópica de $y$ preserva todos los agregados $Z_k$, la Hessiana del
potencial no ve esa dirección y no puede amortiguarla.

\subsection{Densidad predictiva completa}
\label{subsec:densidad-predictiva-completa}

La predicción de congestión se basa en una densidad suavizada:
\begin{equation}
  \rho(x,t)=\sum_i K_h(x-p_i(t)).
  \label{eq:anexo_densidad}
\end{equation}
La primera derivada temporal es
\begin{equation}
  \dot\rho(x,t)
  =
  -\sum_i\nabla K_h(x-p_i)^\top u_i.
  \label{eq:anexo_densidad_dot}
\end{equation}
La segunda derivada completa incluye el término convectivo:
\begin{equation}
  \ddot\rho(x,t)
  =
  \sum_i
  \left[
    u_i^\top\nabla^2K_h(x-p_i)u_i
    -
    \nabla K_h(x-p_i)^\top a_i
  \right].
  \label{eq:anexo_densidad_ddot}
\end{equation}
Si el campo de movimiento es $u_i=-\kappa\nabla_{p_i}U_i$, entonces una expansión útil
para la aceleración es
\begin{equation}
  a_i=
  -\kappa
  \left[
    \nabla_{p_i}^2U_i\,u_i
    +
    \frac{\partial\nabla_{p_i}U_i}{\partial y_i}\dot y_i
    +
    \frac{\partial\nabla_{p_i}U_i}{\partial B_i}\dot B_i
    +\cdots
  \right].
  \label{eq:anexo_aceleracion_campo}
\end{equation}

La densidad futura se aproxima por
\begin{equation}
  \rho^\tau=
  \clip\left[
  \rho+\tau\dot\rho+\frac{\tau^2}{2}\ddot\rho
  \right],
  \label{eq:anexo_rho_tau}
\end{equation}
lo que deforma la velocidad admisible y la eikonal:
\begin{equation}
  v^\tau(x)=v_{\min}+(v_{\max}-v_{\min})e^{-\alpha\rho^\tau(x)},
  \qquad
  \norm{\nabla T_k^\tau(x)}=\frac{1}{v^\tau(x)}.
  \label{eq:anexo_eikonal_predictiva}
\end{equation}

Como $T_k^\tau$ depende de una densidad futura, esta construcción no comparte el mismo
potencial estático que la energía maestra del cuerpo principal. Debe usarse como
extensión predictiva \textbf{[HIPÓTESIS EXTENDIDA]} o como perturbación acotada, no como parte de la
prueba de descenso ideal.

\begin{figure}[H]
\centering
\begin{tikzpicture}[
  block/.style={draw=VIUDark,fill=VIULight,rounded corners=2pt,
    align=center,minimum width=2.85cm,minimum height=0.85cm,font=\small},
  arrow/.style={-{Latex[length=2mm]},thick,color=VIUOrange!90!black}
]
\node[block] (jets) {jets locales\\$p_i,u_i,a_i$};
\node[block, right=0.7cm of jets] (rho) {densidad\\$\rho,\dot\rho,\ddot\rho$};
\node[block, right=0.7cm of rho] (future) {predicción\\$\rho^\tau$};
\node[block, below=0.85cm of future] (vel) {velocidad\\$v^\tau(x)$};
\node[block, left=0.7cm of vel] (eik) {eikonal\\$T_k^\tau$};
\node[block, left=0.7cm of eik] (pay) {payoff\\$\Psi(T_k^\tau)$};
\draw[arrow] (jets)--(rho);
\draw[arrow] (rho)--(future);
\draw[arrow] (future)--(vel);
\draw[arrow] (vel)--(eik);
\draw[arrow] (eik)--(pay);
\end{tikzpicture}
\caption{Densidad predictiva completa: los jets locales generan una eikonal futura que
modula la contribución efectiva antes de entrar a Smith.}
\label{fig:anexo_densidad}
\end{figure}

\subsection{Claridad geodésica}
\label{subsec:claridad-geodesica}

La cota euclidiana para $T$ sólo es válida bajo claridad local. En general debe usarse
distancia geodésica ponderada:
\begin{equation}
  |T(\hat p)-T(p)|\le d_v(\hat p,p).
  \label{eq:anexo_cota_geodesica}
\end{equation}
Si además
\begin{equation}
  \norm{\hat p-p}<\rho_{\mathrm{clear}},
  \label{eq:anexo_condicion_claridad}
\end{equation}
y el segmento local no cruza obstáculos ni cambios de homotopía, entonces
\begin{equation}
  |T(\hat p)-T(p)|
  \le
  \frac{\norm{\hat p-p}}{v_{\min}}.
  \label{eq:anexo_cota_euclidiana_claridad}
\end{equation}
Esta condición es la forma correcta que debe usarse al extender el presupuesto de
percepción a rutas eikonales.
